% Kalkylark internals. Build: make docs   (tectonic)
% Licensed CC BY 4.0.
\documentclass[11pt]{article}
\usepackage{kalkylark}
\usepackage{tikz}
\usetikzlibrary{arrows.meta,positioning,fit}

\tikzset{
  mod/.style={draw=ink, fill=paperbox, rounded corners=2pt, align=center,
    font=\sffamily\small, minimum width=26mm, minimum height=10mm},
  ext/.style={draw=phdim, dashed, rounded corners=2pt, align=center,
    font=\sffamily\small, minimum width=26mm, minimum height=9mm},
  use/.style={-{Stealth[length=2mm]}, thick, color=ink},
  cellnode/.style={draw=ink, fill=paperbox, minimum width=10mm,
    minimum height=6mm, font=\sffamily\small, rounded corners=1pt},
}
\lstdefinestyle{c}{language=C, morekeywords={size_t,ssize_t},
  keywordstyle=\bfseries, commentstyle=\itshape\color{phdim!70!black}}

\title{\sffamily\bfseries\color{ink}Kalkylark internals\\[4pt]
  \Large\mdseries A walk through the source}
\author{\sffamily Set Lonnert and the Kalkylark authors}
\date{\sffamily Version 0.1 · November 2026}

\begin{document}
\maketitle
\thispagestyle{empty}

\begin{abstract}
\noindent Kalkylark is about 2\,800 lines of C99 in seven modules. This
document explains how they work and why they are built the way they are:
a terminal layer on bare \texttt{termios} and escape sequences, a parser
that evaluates while it parses, three recalculation strategies, a number
formatter, an entry editor with pointing, an atomic save, and the testing
setup, which includes a pseudo-terminal harness that also produces the
screens in the documentation. Code excerpts are abridged. The source is
the authority.
\end{abstract}

\vspace{1em}
\begin{center}
\begin{tikzpicture}[node distance=9mm and 8mm, scale=0.9, transform shape]
  \node[mod] (main) {main.c\\\scriptsize options, startup};
  \node[mod, below=of main] (ui) {ui.c\\\scriptsize screen, entry, commands};
  \node[mod, left=of ui] (theme) {theme.c\\\scriptsize phosphor colours};
  \node[mod, below left=of ui, xshift=10mm] (term) {term.c\\\scriptsize keys, VT100 output};
  \node[mod, below right=of ui, xshift=-10mm] (sheet) {sheet.c\\\scriptsize cells, recalc, formats};
  \node[mod, below=of sheet] (formula) {formula.c\\\scriptsize parser, evaluator};
  \node[mod, right=of sheet] (storage) {storage.c\\\scriptsize save, load};
  \node[ext, below=of term] (tty) {terminal\\\scriptsize termios + escapes};
  \node[ext, right=of storage, xshift=-6mm, yshift=-19mm] (file) {.kalk file};
  \draw[use] (main) -- (ui);
  \draw[use] (main) -| (theme);
  \draw[use] (ui) -- (theme);
  \draw[use] (ui) -- (term);
  \draw[use] (ui) -- (sheet);
  \draw[use] (ui) -| (storage);
  \draw[use] (theme) |- (term);
  \draw[use] ([xshift=-3mm]sheet.south) -- ([xshift=-3mm]formula.north);
  \draw[use] ([xshift=3mm]formula.north) -- ([xshift=3mm]sheet.south);
  \draw[use] (storage) -- (sheet);
  \draw[use] (term) -- (tty);
  \draw[use] (storage) |- (file);
\end{tikzpicture}

{\small\sffamily\color{ink}Modules and who calls whom. \texttt{sheet.c}
and \texttt{formula.c} call each other: the sheet evaluates formulas, and
formulas read cells.}
\end{center}

\newpage
\tableofcontents
\newpage

% ======================================================================
\section{Principles}

\begin{description}[style=nextline, leftmargin=1.5em]
\item[Portable by construction.] C99 plus POSIX.1-2008, and nothing else:
  no curses, no libraries beyond \texttt{libm}. The code builds without a
  warning with \texttt{-std=c99 -pedantic}, \texttt{-Wall -Wextra} and
  \texttt{-Wconversion -Wshadow}.
\item[The engine knows nothing about the screen.] \texttt{sheet.c},
  \texttt{formula.c} and \texttt{storage.c} form the \emph{core}. They
  never include a terminal header, and the tests link only the core.
\item[Fixed-size where the original was.] VisiCalc's sheet had 63 columns
  and 254 rows. A static array of 16\,002 cells is small (about 400\,KB),
  needs no growth logic, and makes coordinates plain array indices.
\item[Faithful first, then optional.] VisiCalc's behaviour is the default
  (left-to-right formulas, one-pass recalculation). Later improvements,
  such as natural order, are options.
\end{description}

% ======================================================================
\section{The terminal layer: \texttt{term.c}}

Curses solves terminal differences with the \texttt{terminfo} database.
Today every terminal emulator understands the VT100/ECMA-48 escape
sequences, so Kalkylark writes them directly and needs only \texttt{termios}
for input.

\subsection{Raw mode}

By default a terminal is in \emph{canonical} mode: the kernel collects a
whole line, echoes it, and turns \kbd{Ctrl-C} into a signal. A spreadsheet
wants every key as it is pressed:

\begin{lstlisting}[style=c]
raw = orig_tio;
raw.c_iflag &= ~(tcflag_t)(BRKINT | ICRNL | INPCK | ISTRIP | IXON);
raw.c_oflag &= ~(tcflag_t)OPOST;     /* we position everything ourselves */
raw.c_cflag |= CS8;
raw.c_lflag &= ~(tcflag_t)(ECHO | ICANON | IEXTEN | ISIG);
raw.c_cc[VMIN] = 1;                  /* read() returns after one byte  */
raw.c_cc[VTIME] = 0;
tcsetattr(STDIN_FILENO, TCSAFLUSH, &raw);
\end{lstlisting}

Turning off \texttt{ISIG} means \kbd{Ctrl-C} arrives as byte 3 and is
handled as a key, and turning off \texttt{IXON} frees \kbd{Ctrl-S} and
\kbd{Ctrl-Q}. The original settings are saved and restored by
\texttt{term\_done()}. That function is registered with \texttt{atexit},
called on \texttt{SIGTERM} and \texttt{SIGHUP}, and safe to call twice. A
crash must never leave the user's shell in raw mode. On entry Kalkylark
also switches to the \emph{alternate screen} (\texttt{ESC[?1049h}), so the
shell's scrollback is untouched when it quits.

\subsection{Decoding keys}

Arrow and function keys arrive as escape sequences: \texttt{ESC [ A} for
\up, \texttt{ESC [ 5 \textasciitilde} for \kbd{PgUp}, and so on. The
difficulty is the \kbd{Esc} key itself, which sends a lone \texttt{ESC}.
The decoder waits up to 50\,ms for a following byte, using
\texttt{poll()}. If none comes, it was \kbd{Esc}:

\begin{lstlisting}[style=c]
c = read_byte(ESC_TIMEOUT_MS);       /* -1 on timeout */
if (c < 0)
    return KEY_ESC;
if (c == '[' || c == 'O')
    return decode_csi();             /* read up to a final byte 0x40..0x7e */
\end{lstlisting}

\texttt{decode\_csi()} collects parameter bytes until a \emph{final byte}
in the range \texttt{0x40}--\texttt{0x7e}, as ECMA-48 defines it. As a
result, modified keys such as \texttt{ESC[1;5C} (\kbd{Ctrl}+\rgt) are
recognised as plain arrows instead of leaving stray characters behind.

\subsection{Output and redrawing}

All output goes into a 64\,KB buffer and leaves in a single
\texttt{write()} per frame. The whole screen is redrawn after every key,
which costs about 3\,KB at 80\,\(\times\)\,24 and avoids any bookkeeping
about what changed. Because each line is overwritten in place and ended
with \texttt{ESC[K} (erase to end of line), the screen never flickers.

\subsection{Window size}

\texttt{TIOCGWINSZ} gives the size. \texttt{SIGWINCH} is installed
\emph{without} \texttt{SA\_RESTART}, so it interrupts the blocking
\texttt{read()}. The key reader then returns \texttt{KEY\_RESIZE}, and the
next frame is drawn at the new size.

\begin{note}[Portability]
\texttt{SIGWINCH} and \texttt{TIOCGWINSZ} are on every Unix but only
entered POSIX in 2024, and macOS hides them in strict POSIX mode.
\texttt{term.c} defines \texttt{\_DARWIN\_C\_SOURCE} on Apple systems and
guards both with \texttt{\#ifdef}. A system without them falls back to
\texttt{\$LINES}/\texttt{\$COLUMNS}, or 80\,\(\times\)\,24.
\end{note}

% ======================================================================
\section{Colours: \texttt{theme.c}}

The screen uses five \emph{roles}: normal, dim, bright, border and cursor.
A theme turns each role into an SGR sequence (\texttt{ESC[...m}) once, at
start-up, at the colour depth the terminal supports: 16 colours, the
256-colour palette, or 24-bit \texttt{38;2;r;g;b}.

A phosphor has a single hue at three intensities on a near-black
background. The border and the cell cursor are not separate colours. They
are \emph{reverse video} (SGR~7) of the dim and bright intensities, just as
on a monochrome monitor:

\begin{lstlisting}[style=c]
snprintf(seq[ROLE_NORMAL], ..., "\033[0;%sm",   norm);
snprintf(seq[ROLE_BORDER], ..., "\033[0;7;%sm", dim);
snprintf(seq[ROLE_CURSOR], ..., "\033[0;7;%sm", bright);
\end{lstlisting}

The \texttt{plain} theme uses only bold and reverse video, so it works on
any terminal and respects the user's own colour scheme. \texttt{NO\_COLOR}
selects it unless \texttt{-p} is given.

% ======================================================================
\section{Cells: \texttt{sheet.c}}

\begin{lstlisting}[style=c]
struct val  { enum val_kind kind; double num; };     /* NUM, NA, ERR */
struct cell {
    enum cell_type type;    /* EMPTY, LABEL, VALUE */
    char fmt;               /* display format, 0 = global */
    char *text;             /* label text, or the formula as typed */
    struct val v;           /* last computed value */
};
static struct cell cells[SHEET_ROWS][SHEET_COLS];
\end{lstlisting}

A value cell keeps its formula as \emph{text}, uppercased, plus the last
computed value. There is no compiled form. Re-parsing at every
recalculation is fast enough for 16\,002 cells, and it keeps one source of
truth.

\texttt{NA} and \texttt{ERR} are kinds of value, not numbers. Arithmetic
propagates them, with \texttt{ERR} winning, and any non-finite result
(division by zero, \texttt{@SQRT(-1)}, overflow) becomes \texttt{ERR}.

Whether an entry is a label or a value is decided by its first character,
as in VisiCalc: a digit or one of \texttt{+ - . ( @} makes a value, and a
leading \texttt{"} forces a label.

% ======================================================================
\section{Formulas: \texttt{formula.c}}

\subsection{Evaluate while parsing}

The parser is recursive descent, and it computes values as it recognises
them. There is no syntax tree:

\begin{lstlisting}
expr    := unary { op unary }        op: + - * / ^   (left to right)
unary   := '-' unary | '+' unary | primary
primary := number | ref | '(' expr ')' | '@' function
\end{lstlisting}

VisiCalc's lack of precedence makes the expression rule a single loop.
That small loop is the whole reason \texttt{2+3*4} is 20:

\begin{lstlisting}[style=c]
static struct val parse_expr(struct parser *p)
{
    struct val a = parse_unary(p);
    for (;;) {
        skip(p);
        int op = *p->s;
        if (op == '\0' || strchr("+-*/^", op) == NULL || p->err)
            return a;
        p->s++;
        a = binop(a, op, parse_unary(p));
    }
}
\end{lstlisting}

The usual way to add precedence is one function per level
(\texttt{term}, \texttt{factor}, \texttt{power}). A \texttt{/G} option
for it would be a small, self-contained change.

\subsection{One parser, three uses}

The same code runs in three modes, chosen by how \texttt{struct parser} is
set up:

\begin{center}
\begin{tabularx}{\linewidth}{@{}lX@{}}
\toprule
\texttt{formula\_check()} & syntax only: is the entry acceptable? Returns a
  message and the column of the first error. \\
\texttt{formula\_eval()} & the value, read from the cells' current values. \\
\texttt{formula\_refs()} & every cell the formula may read, reported through
  a callback. Natural-order recalculation builds its graph from this. \\
\bottomrule
\end{tabularx}
\end{center}

Errors are recorded once (\texttt{fail()} keeps the first) and parsing
continues harmlessly. That keeps the code free of error checks at every
call site.

\subsection{References and ranges}

A reference is one to three letters followed by one to four digits, and
nothing alphanumeric after it. That is how \texttt{SUM} without an
\texttt{@} is rejected rather than read as a cell. A range is a reference,
then \texttt{.}, \texttt{..}, \texttt{...} or \texttt{:}, then another
reference. Reversed corners are normalised. In \texttt{formula\_refs} mode
a range reports every cell in it, and \texttt{@LOOKUP} also reports the
column beside its range, because which cell the result comes from depends
on the data.

% ======================================================================
\section{Recalculation}

\subsection{The sweep}

\begin{lstlisting}[style=c]
for (c = 0; c < SHEET_COLS; c++)           /* /GOR swaps the loops */
    for (r = 0; r < SHEET_ROWS; r++)
        if (cells[r][c].type == CELL_VALUE)
            cells[r][c].v = formula_eval(cells[r][c].text);
\end{lstlisting}

\subsection{Natural order with Kahn's algorithm}

Lotus 1-2-3 is usually described with a recursive depth-first walk. A
chain of 16\,002 formulas would then mean 16\,002 nested calls, each with
the parser's frames on top, which is enough to overflow an 8\,MB stack.
Kalkylark uses Kahn's algorithm instead, which is iterative:

\begin{enumerate}[label=\arabic*.]
\item \textbf{Edges.} For each formula cell \(x\), run
  \texttt{formula\_refs} and record an edge \(p \to x\) for every formula
  cell \(p\) that \(x\) reads. Cells holding plain numbers or labels need
  no ordering and are skipped.
\item \textbf{Compressed adjacency.} Count the dependents of each cell,
  take prefix sums, and fill one array: the dependents of \(p\) are
  \texttt{adj[first[p]..first[p+1])}. The in-degree of each cell is its
  number of precedents.
\item \textbf{Queue.} Put every formula with in-degree 0 in a queue, in
  column order so the result is deterministic. Take one, compute it,
  decrement each dependent's in-degree, and queue any that reach 0.
\item \textbf{Leftovers are cycles.} Anything never queued is on a cycle
  or downstream of one. Compute it in column order, VisiCalc style, and
  set \texttt{sheet\_circular}.
\end{enumerate}

\begin{lstlisting}[style=c]
while (head < tail) {
    x = queue[head++];
    evaluate(x);
    done[x] = 1;
    for (i = first[x]; i < first[x + 1]; i++)
        if (--indeg[adj[i]] == 0)
            queue[tail++] = adj[i];
}
\end{lstlisting}

The graph is rebuilt on every recalculation, which costs a second parse
of each formula. That is simple and always correct. The test that fills
all 16\,002 cells with a single chain runs in a fraction of a second. If
memory runs out, the code falls back to column order instead of failing.

\begin{center}
\begin{tikzpicture}[node distance=6mm and 14mm]
  \node[cellnode] (a20) {A20};
  \node[cellnode, right=of a20] (b5) {B5};
  \node[cellnode, right=of b5] (c9) {C9};
  \node[cellnode, right=of c9] (a1) {A1};
  \draw[use] (a20) -- (b5);
  \draw[use] (b5) -- (c9);
  \draw[use] (c9) -- (a1);
  \node[below=1mm of a20, font=\sffamily\scriptsize, color=ink] {3};
  \node[below=1mm of b5,  font=\sffamily\scriptsize, color=ink] {+A20};
  \node[below=1mm of c9,  font=\sffamily\scriptsize, color=ink] {+B5*10};
  \node[below=1mm of a1,  font=\sffamily\scriptsize, color=ink] {+C9};
\end{tikzpicture}

{\small\sffamily\color{ink}A zig-zag chain from the tests: column order
needs four passes, Kahn's algorithm one.}
\end{center}

% ======================================================================
\section{Showing numbers}

\texttt{format\_number()} implements the general format: the most precise
form that fits in the column.

\begin{enumerate}[label=\arabic*.]
\item Try fixed point with 10, 9, \dots, 0 decimals, dropping trailing
  zeros, and take the first that fits. Never accept a result that rounds a
  non-zero number to \texttt{0}.
\item Otherwise use scientific notation with as many digits as fit, in a
  compact form: \texttt{1.2346e8}, not \texttt{1.2346e+08}.
\item Otherwise fill the column with \texttt{>}.
\end{enumerate}

The other formats (\texttt{I}, \texttt{\$}, \texttt{*}) are a few lines
each. \texttt{-0} and \texttt{-0.00} are normalised to \texttt{0} and
\texttt{0.00}.

% ======================================================================
\section{The screen and the keyboard: \texttt{ui.c}}

\subsection{One function draws everything}

\texttt{draw()} rebuilds the status, prompt and entry lines, the border,
and the visible part of the grid from the current state. It first scrolls
so that the highlighted cell is visible. That cell is the cursor, or the
pointer while pointing. Every interaction is a loop of \emph{draw, read a
key, change state}. Nested prompts (\texttt{ask\_key}, \texttt{ask\_line})
run their own loops and redraw on every key, including on window resize.

\subsection{Entry and pointing}

\texttt{cell\_entry()} is a small state machine around an edit buffer:

\begin{center}
\begin{tikzpicture}[node distance=8mm and 26mm,
    st/.style={draw=ink, fill=paperbox, rounded corners=6pt,
      font=\sffamily\small, minimum width=22mm, minimum height=8mm}]
  \node[st] (typing) {typing};
  \node[st, right=of typing] (pointing) {pointing};
  \node[st, below=12mm of typing, xshift=24mm] (stored) {stored};
  \draw[use] (typing) to[bend left=15] node[above, font=\sffamily\scriptsize, align=center]
    {arrow after\\\texttt{+-*/\^{}(,.}} (pointing);
  \draw[use] (pointing) to[bend left=15] node[below, font=\sffamily\scriptsize]
    {any other key} (typing);
  \draw[use] (pointing) to[loop right] node[right, font=\sffamily\scriptsize, align=left]
    {arrow: move pointer,\\rewrite reference} (pointing);
  \draw[use] (typing) -- node[left, font=\sffamily\scriptsize] {Enter / arrow} (stored);
  \draw[use] (pointing) -- node[right, font=\sffamily\scriptsize] {Enter} (stored);
\end{tikzpicture}
\end{center}

While pointing, the buffer position where the reference starts is
remembered (\texttt{ref\_at}). Each arrow truncates the buffer there and
appends the pointer's new coordinate. Typing \texttt{.} while pointing
appends \texttt{...} and leaves pointing mode, so the next arrow starts the
second corner of a range from where the pointer was. If the formula is
rejected, the error stays on the prompt line and the entry remains open for
editing.

\subsection{Prompts with defaults}

\texttt{ask\_line()} can start from a suggested value, such as the current
file name or column width. Enter accepts it, the first typed character
replaces it, and Backspace edits it. The column-width prompt originally
appended to the suggestion, so typing \texttt{12} over \texttt{9} gave
\texttt{912}. The end-to-end tests caught that.

% ======================================================================
\section{Files: \texttt{storage.c}}

A sheet file is a list of records that could almost be typed in: global
settings (\texttt{/GO}, \texttt{/GC}, \texttt{/GF}) and one line per
non-empty cell, \texttt{>B3:/F\$:@SUM(B1...B2)}. Labels are always stored
with their \texttt{"}, so loading can feed every cell through
\texttt{sheet\_enter()}, the same function the keyboard uses. The loader
then has no parsing of its own to get wrong, and a bad formula in a file is
reported with its line number.

Saving is atomic: write \emph{name}\texttt{.tmp}, check \texttt{ferror}
and \texttt{fclose}, then \texttt{rename()} it over the target. POSIX
guarantees that \texttt{rename} replaces the file in one step, so a full
disk or a crash leaves the old file intact.

% ======================================================================
\section{Testing}

\begin{description}[style=nextline, leftmargin=1.5em]
\item[\texttt{make check}] builds \texttt{test\_core}, which links only the
  core and checks about a hundred cases: arithmetic order, references,
  every function, \texttt{NA}/\texttt{ERROR} propagation, syntax errors,
  number formatting, every display format, a save/load round trip
  (including a label containing \texttt{:} and \texttt{"}), and natural
  order with zig-zag chains, ranges, \texttt{@LOOKUP}, cycles, and the
  longest possible chain (all 16\,002 cells). It then loads every sheet in
  \texttt{examples/} and fails if any cell shows \texttt{ERROR}. That check
  found a real mistake in the loan example.
\item[\texttt{tools/screens.py}] starts the real program in a
  pseudo-terminal, sends keystrokes, and interprets the escape sequences it
  writes, cursor positioning, erasing and SGR colours, into a character
  grid with attributes. It writes LaTeX for the documentation, so every
  screen in the tutorial is a test that the UI still behaves as described.
  One subtlety: an escape sequence can be split across two reads, so the
  interpreter keeps any incomplete tail for the next read.
\item[CI] builds and tests on Linux, macOS and FreeBSD, and builds the
  PDFs, on every push.
\end{description}

% ======================================================================
\section{Where to extend}

\begin{center}
\begin{tabularx}{\linewidth}{@{}lX@{}}
\toprule
\textbf{Feature} & \textbf{Where and how} \\
\midrule
\texttt{/R} replicate & \texttt{ui.c} and a reference rewriter in
  \texttt{formula.c}: copy a formula, shifting relative references (VisiCalc
  asked ``No change / Relative'' for each one). \\
\texttt{/I}, \texttt{/D} insert, delete & move rows/columns in
  \texttt{sheet.c} and rewrite every reference that points past the
  change, using the same rewriter. \\
\texttt{/T} titles, \texttt{/W} windows & \texttt{ui.c}: frozen rows or
  columns, and a split screen with two viewports. \\
Column widths & per-column array in \texttt{sheet.c}. The file format
  needs one more record type. \\
Minimal recalculation & keep the graph from \texttt{recalc\_natural}
  between recalculations, add dirty flags, and recompute only what is
  reachable from the edited cell. \\
Precedence option & a second expression grammar in \texttt{formula.c},
  selected by a \texttt{/G} setting saved in the file. \\
\bottomrule
\end{tabularx}
\end{center}

\vfill
{\small This document is licensed CC BY 4.0. Kalkylark is free software
under the BSD 2-clause licence.}

\end{document}
